\section{Setting}
\label{sec:setting}

Unless otherwise specified, finite-dimensional quantities depend
on $d$, with this dependence suppressed when unambiguous.
We write $\mathscr P_2(E)$ for the Borel probability measures on
$E\subseteq\mathbb R^k$ with finite second moment, $W_2$ for the
$2$-Wasserstein distance, and $\delta_x$ for the Dirac measure at $x$.
See Appendix~\ref{subsec:wasserstein} for further details on Wasserstein spaces.

\subsection{Model}

\begin{definition}[Binary Gaussian Mixture Model]
\label{def:gaussian-mixture-model}
Let $d\ge1$ be the dimension, $n\ge1$ the sample size,
$\mu\in\mathbb R^d$ the signal vector, and
$Y\in\{-1,+1\}^n$, $\Delta\in\{0,1\}^n$
the vectors of true labels and supervision indicators,
respectively. These quantities are deterministic.
Fix a noise scale $\sigma>0$. The feature vectors are
generated according to
\begin{equation}
    X_i
    =
    \frac{Y_i\mu}{\sqrt d}+\sigma U_i,
    \qquad
    U_i\stackrel{\mathrm{i.i.d.}}{\sim}
    \mathcal N(0_d,I_d),
    \qquad i\in[\![1,n]\!].
    \label{eq:gaussian-mixture-model}
\end{equation}
Thus, the observations are independent Gaussian vectors,
with class memberships determined by the fixed labels.
Only labels with $\Delta_i=1$ are observed, giving the
partially labeled dataset
\[
    \mathcal D
    :=
    \{(X_i,\Delta_iY_i):i\in[\![1,n]\!]\}.
\]
\end{definition}
For later use, define the labeled and unlabeled index sets
\[
    I_{\mathcal L}
    :=
    \{i\in[\![1,n]\!]:\Delta_i=1\},
    \qquad
    I_{\mathcal U}
    :=
    \{i\in[\![1,n]\!]:\Delta_i=0\},
\]
with cardinalities
$N:=\#I_{\mathcal L}$ and $M:=\#I_{\mathcal U}$.
Write $\delta := n / d$ for the sample-to-dimension ratio, $\rho := N/n$ for the labeled proportion,
and $p=\#\{i: Y_i=+1\}/n$ for the positive-class proportion.
We also define the empirical signal-coordinate and
sample-type laws
\begin{equation}
    \nu_\mu
    :=
    \frac1d\sum_{k=1}^d\delta_{\mu_k},
    \qquad
    \nu_{(Y,\Delta)}
    :=
    \frac1n\sum_{i=1}^n\delta_{(Y_i,\Delta_i)}.
    \label{eq:empirical-environment-laws}
\end{equation}
General empirical-measure notation is introduced in
Section~\ref{subsec:empirical-measures}.
Classical asymptotic statistics typically considers a fixed
parameter dimension as the sample size grows.
When the dimension is comparable to the sample size,
classical approximations may fail~\cite{sur2019modern}.
Nevertheless, aggregate quantities can exhibit predictable
limiting behavior even when individual coordinates remain random.
This motivates the \emph{proportional asymptotic regime}, in which
\[
    n,d\longrightarrow\infty,
    \qquad
    \frac nd\longrightarrow\bar\delta\in(0,\infty).
\]
More precisely, the proportional asymptotic regime considers
a sequence of statistical models indexed by the dimension $d$,
rather than a single fixed-dimensional model.
The sample size $n=n(d)$, signal vector $\mu=\mu(d)$,
labels $Y=Y(d)$, and supervision indicators
$\Delta=\Delta(d)$ all depend on $d$.
By imposing suitable conditions on their asymptotic behavior,
we can obtain precise limiting descriptions as $d\to\infty$.
We suppress this dependence on $d$ in the notation and
formalize these conditions in the following assumption.
\begin{assumption}[Asymptotic Regime]
\label{assumption:asymptotic-regime}
As $d\to\infty$, assume that
\begin{equation}
    \delta\longrightarrow\bar\delta,
    \qquad
    \nu_\mu\xrightarrow{W_2}\bar\nu_\mu,
    \qquad
    \nu_{(Y,\Delta)}
    \xrightarrow{W_2}\bar\nu_{(Y,\Delta)},
    \label{eq:finite-dim-limits}
\end{equation}
with $\bar \delta\in(0,\infty)$.
In particular, the positive-class and labeled proportions
converge:
\begin{equation}
    p\longrightarrow\bar p,
    \qquad
    \rho\longrightarrow\bar\rho.
    \label{eq:finite-dim-limits-2}
\end{equation}
We assume $\bar p,\bar\rho\in(0,1)$.
Consequently, $N,M\to\infty$, and both are positive
for all sufficiently large $d$.
\end{assumption}
In what follows, we restrict attention to dimensions
for which $N,M\ge1$.

\subsection{Self-Training Algorithm}

We now formalize the self-training algorithm introduced in
Section~\ref{sec:introduction}.
For $\theta=(b,w)\in\mathbb R^{d+1}$, define the
augmented feature vector and score function
\[
    \breve X_i
    :=
    \begin{bmatrix}
        1\\
        X_i/\sqrt d
    \end{bmatrix},
    \qquad
    r_i(\theta)
    :=
    \breve X_i^\top\theta
    =
    b+\frac{X_i^\top w}{\sqrt d},
    \qquad
    i \in [\![1,n]\!].
\]
Let $\ell:\{-1,+1\}\times\mathbb R\to\mathbb R$
be a differentiable loss function,
$j:\mathbb R\to\mathbb R$ a differentiable penalty function, and let $\kappa := (\kappa_-,\kappa_+)$ be selection thresholds, with $\kappa_- < 0 < \kappa_+$.
We write
\[
    \ell'(y,r)
    :=
    \frac{\partial}{\partial r}\ell(y,r),
    \qquad
    J(w):=\sum_{k=1}^d j(w_k),
    \qquad 
    S_\kappa(r) := \mathbf1_{ \{ r \le \kappa_-\} } + \mathbf1_{ \{ r \ge \kappa_+\} }.
\]
For a teacher parameter $\vartheta \in \mathbb{R}^{d+1}$ and $j \in I_{\mathcal U}$, the associated pseudo-label is $\widehat Y_j(\vartheta):=\operatorname{sign}(r_j(\vartheta))$, and
the confidence-selection step in \eqref{eq:pl-dataset}
can be expressed through the selection indicator $S_j(\vartheta) := S_\kappa(r_j(\vartheta))$. 
The indicator $S_j(\vartheta)$ ensures that unselected
observations do not contribute to the training objective.
We use the convention $\operatorname{sign}(0)=+1$.

Fix dimension-independent coefficients $\lambda,\pi\ge0$.
The self-training objective, with teacher parameter
$\vartheta$ and student parameter $\theta$, is
\begin{align}
    L_{\mathrm{ST}}(\theta;\vartheta)
    :=
    &\frac1N\sum_{i\in I_{\mathcal L}}
        \ell\bigl(Y_i,r_i(\theta)\bigr)
    +
    \frac{\pi}{M}\sum_{j\in I_{\mathcal U}}
        S_j(\vartheta)
        \ell\bigl(\widehat Y_j(\vartheta),r_j(\theta)\bigr)
    +
    \frac{\lambda}{d}J(w).
    \label{eq:ST-objective}
\end{align}
At iteration $t$, this is the objective
$L(\theta;\mathcal D_L,\mathcal P_t)$ introduced in
\eqref{eq:pl-gd-step}, with
$\vartheta=\theta^t$ held fixed.

We study the procedure in which the teacher is refreshed
after every gradient step. Starting from
$\theta^0=(b^{\mathrm{init}},w^{\mathrm{init}})$,
the iteration is
\begin{equation}
    \theta^{t+1}
    :=
    \theta^t-\eta P_d
    \nabla_1L_{\mathrm{ST}}(\theta^t;\theta^t),
    \qquad t\ge0,
    \label{eq:ST-recursion}
\end{equation}
where $\nabla_1$ denotes differentiation with respect
to the student parameter only, $\eta>0$ is the step size,
and
\[
    P_d
    :=
    \begin{bmatrix}
        1 & 0_d^\top\\
        0_d & dI_d
    \end{bmatrix}
    \in\mathbb R^{(d+1)\times(d+1)}.
\]
This preconditioning gives a nontrivial scaling of the
parameter updates in the proportional regime.
The explicit updates are summarized below.

\begin{algorithm}[Self-Training]
\label{alg:self-training}
Fix $d\ge1$ with $N,M\ge1$.
\begin{itemize}[wide]
    \item \textbf{Input:}
    Dataset $\mathcal D$, initialization
    $(b^{\mathrm{init}},w^{\mathrm{init}})$,
    number of iterations $T\in\mathbb N$,
    and hyperparameters
    $\eta$, $\lambda,\pi\ge0$ and $\kappa_-<0<\kappa_+$.

    \item \textbf{Initialization:}
    Set
    \[
        b^0:=b^{\mathrm{init}},
        \qquad
        w^0:=w^{\mathrm{init}}.
    \]

    \item \textbf{Iterations:}
    For $t=0,\ldots,T-1$:

    Compute the scores, and the pseudo-labels and selection indicators:
    \begin{align*}
        &r_i^t
        :=
        b^t+\frac{X_i^\top w^t}{\sqrt d},
         & &i\in[\![1,n]\!], \\
        \widehat Y_j^t :=\operatorname{sign}(r_j^t), &
        \qquad
        S_j^t
        :=
        \mathbf 1_{\{r_j^t\le\kappa_-\}}
        +
        \mathbf 1_{\{r_j^t\ge\kappa_+\}}
        , &&j\in I_{\mathcal U}
        .
    \end{align*}
    
    Update the intercept and weights:
    \begin{align*}
        b^{t+1}
        &:=
        b^t-\eta\Biggl[
            \frac1N\sum_{i\in I_{\mathcal L}}
                \ell'(Y_i,r_i^t)
            +
            \frac{\pi}{M}\sum_{j\in I_{\mathcal U}}
                S_j^t\ell'(\widehat Y_j^t,r_j^t)
        \Biggr],\\
        w^{t+1}
        &:=
        w^t-\eta d\Biggl[
            \frac1N\sum_{i\in I_{\mathcal L}}
                \ell'(Y_i,r_i^t)\frac{X_i}{\sqrt d}
            +
            \frac{\pi}{M}\sum_{j\in I_{\mathcal U}}
                S_j^t\ell'(\widehat Y_j^t,r_j^t)
                \frac{X_j}{\sqrt d}
            +
            \frac{\lambda}{d}\nabla J(w^t)
        \Biggr].
    \end{align*}

    \item \textbf{Output:}
    The classifier $f_{\theta^T}$,
    where $\theta^T=(b^T,w^T)$.
\end{itemize}
\end{algorithm}

The assumptions required for our main results differ and
are introduced alongside the corresponding theorems.
Nevertheless, all conditions imposed on the loss and penalty
are satisfied by the logistic loss and ridge penalty
\begin{equation}
    \ell(y,r):=\log(1+e^{-yr}),
    \qquad
    j(u):=\frac{u^2}{2}.
    \label{eq:logistic-loss-ridge-penalty}
\end{equation}
